\documentclass[a4paper,12pt]{article}
\usepackage[italian]{babel}
\usepackage[utf8]{inputenc}
\usepackage[T1]{fontenc}
\usepackage{geometry}
\usepackage{setspace}
\usepackage{CJKutf8}
\usepackage{fancyhdr}
\geometry{margin=2.5cm}
\onehalfspacing
\pagestyle{fancy}
\fancyhf{}
\renewcommand{\headrulewidth}{0pt}
\fancyhead[C]{\begin{CJK}{UTF8}{min}これは公式試験ですので、解答を提供しないでください。\end{CJK}}
\fancyfoot[C]{\begin{CJK}{UTF8}{min}これは公式試験ですので、解答を提供しないでください。\end{CJK}}

\title{Verifica scritta di informatica - Classe 4 AIQ }
\author{Prof. Fedeli Massimo}
\date{25/03/2026}

\begin{document}
	\maketitle
	
	\textbf{Durata:} 2 ore \\
	\textbf{Obiettivo:} scrivere interrogazioni SQL (anche nidificate) e formulare interrogazioni in algebra relazionale.
	
	\section*{Schema logico}
	
	\begin{flushleft}
		STUDENTI(id\_studente PK, cognome, nome, email, data\_iscrizione) \\
		CORSI(id\_corso PK, titolo, categoria, livello, costo, attivo) \\
		DOCENTI(id\_docente PK, cognome, nome, specializzazione) \\
		LEZIONI(id\_lezione PK, id\_corso FK$\rightarrow$CORSI, id\_docente FK$\rightarrow$DOCENTI, data\_lezione, durata) \\
		ISCRIZIONI(id\_iscrizione PK, id\_studente FK$\rightarrow$STUDENTI, id\_corso FK$\rightarrow$CORSI, data\_iscrizione, stato) \\
		PAGAMENTI(id\_pagamento PK, id\_iscrizione FK$\rightarrow$ISCRIZIONI, data\_pagamento, importo, metodo, esito) \\
		RECENSIONI(id\_recensione PK, id\_studente FK$\rightarrow$STUDENTI, id\_corso FK$\rightarrow$CORSI, voto, commento, data\_recensione)
	\end{flushleft}
	
	\section*{Parte A -- Query SQL - 7.5 pt}
	
	Scrivere le seguenti interrogazioni in linguaggio SQL.
	
	\begin{enumerate}
		\item Mostrare tutte le iscrizioni dello studente con id\_studente = 5 indicando studente, corso e data di iscrizione, ordinate per data decrescente.
		
		\item Elencare i corsi attivi della categoria ``Informatica'' con costo maggiore di 100 euro.
		
		\item Elencare cognome, nome e titolo del corso degli studenti iscritti ai corsi di livello ``Base''.
		
		\item Elencare i corsi che hanno almeno 10 iscritti.
		
		\item Elencare le iscrizioni che non hanno un pagamento con esito positivo.
		
		\item Per ogni corso calcolare il numero totale di studenti iscritti, mostrando solo i corsi con almeno 5 iscritti.
		
		\item Per ogni studente calcolare la spesa totale sostenuta considerando solo i pagamenti con esito positivo, ordinando per spesa decrescente.
		
		\item Per ogni corso calcolare il voto medio delle recensioni, mostrando titolo del corso e media voti.
		
		\item Trovare gli studenti iscritti a tutti i corsi attivi.
		
		\item Per ogni corso calcolare l'incasso totale (somma dei pagamenti con esito positivo). Mostrare titolo del corso e incasso totale.
	\end{enumerate}
	
	\section*{Parte B -- Algebra Relazionale 2.5 pt}
	
	Scrivere le seguenti interrogazioni in algebra relazionale.
	
	\begin{enumerate}
		\item Nomi dei corsi attivi di livello ``Avanzato''.
		\item Cognome e nome degli studenti iscritti nel 2026.
		\item Titoli dei corsi che compaiono in almeno una iscrizione.
		\item Coppie (studente, corso) relative alle iscrizioni del mese di gennaio 2026.
	\end{enumerate}
	
\end{document}